%------------------------------
\begin{myslide}{c}{Automatic Itemize Animations with \lstinline{[<+->]}}

\begin{itemize}[<+->]

\item To see the animations, ensure that \lstinline{handout} in \lstinline{slides.tex} is removed from the \lstinline{documentclass} options

\item \lipsum[1][1-2]

\item \lipsum[1][3]

\item Here we have a \myterm{Term in an animated slide}

\begin{itemize}

\item \lipsum[1][3]

\item \lipsum[1][4-5]

\end{itemize}

\item \lipsum[1][4-5]

\item \lipsum[1][6]

\end{itemize}
\annotation{
This slide demonstrates automatic animations using the [<+->] option in itemize environments. Each item appears sequentially when advancing slides. Note that animations only work when the handout option is removed from the documentclass. Nested lists also animate properly.
}
\end{myslide}

%------------------------------
\begin{myslide}{c}{Animations Custom Ordering}

\begin{columns}

\begin{column}{0.65\textwidth}

\begin{itemize}

\uncover<1->{%
\item \lipsum[1][1-2]
}

\uncover<4->{%
\item \lipsum[1][3]
}

\uncover<5->{%
\item \lipsum[1][4-5]
}

\uncover<2->{%
\item \lipsum[1][6]
}

\end{itemize}

\end{column}

\uncover<3->{%
\myfigcol{\mygraphicspath/csu.pdf}{0.25}
}

\end{columns}
\annotation{
This slide demonstrates custom animation ordering using the uncover command with specific frame numbers. Items can appear in any order you specify (e.g., item 1, then item 4, then item 5, then item 2), and figures can also be animated. This gives precise control over the presentation sequence.
}
\end{myslide}

%------------------------------
\begin{myslidefragile}{c}{Special Call-Out Boxes (2) --- Animated}

\begin{columns}

\begin{column}{1\textwidth}

\begin{itemize}[<+->]

\itembox
\begin{myremark}
\lipsum[1][1]
\end{myremark}

\item \lipsum[1][6]%

\itembox
\begin{myimportant}
\lipsum[1][1-2]
\end{myimportant}

\itembox
\begin{myfuture}
\lipsum[1][1]
\end{myfuture}

\item \lipsum[1][6]

\begin{itemize}

\item Note use of \lstinline{\unskip} for better spacing with boxes after nested lists

\end{itemize}

\unskip

\itembox
\begin{myquestion}
\lipsum[1][2-3]
\end{myquestion}

\end{itemize}

\end{column}

\end{columns}
\annotation{
This slide shows that call-out boxes can also be animated using the [<+->] option. Each box and list item appears sequentially, creating a progressive reveal effect. This is useful for building up complex information step by step.
}
\end{myslidefragile}

%------------------------------
\begin{myslide}{c}{\lstinline{myfigcol} Command --- Animated}

\begin{columns}

\uncover<+->{%
\begin{column}{0.7\textwidth}
\lipsum[66]
\end{column}%
}

\uncover<+->{%
\myfigcol{\mygraphicspath/csu.pdf}{0.15}%
}

\end{columns}

\begin{columns}[t]

\uncover<+->{%
\myfigcol{\mygraphicspath/csu.pdf}{0.25}%
}

\uncover<+->{%
\myfigcol{\mygraphicspath/csu.pdf}{0.20}%
}

\uncover<+->{%
\myfigcol{\mygraphicspath/csu.pdf}{0.15}%
}

\end{columns}
\annotation{
This slide demonstrates animating figures using the uncover command with the <+-> syntax, which automatically increments the frame number. Text and figures can appear sequentially, creating a dynamic presentation flow.
}
\end{myslide}

%------------------------------
\begin{myslidefragile}{c}{Animated Media in WebPressive}

\begin{columns}[t]

\begin{column}{0.45\textwidth}
\centering
\animstill{27mm}{videos/latex_talk_poster.png}{videos/latex_talk.mp4}{https://www.youtube.com/watch?v=N17Od3rY0bA}

\medskip
MP4 video
\end{column}

\begin{column}{0.45\textwidth}
\centering
\animstill{27mm}{videos/latex_talk_poster.png}{videos/latex_talk.gif}{https://www.youtube.com/watch?v=N17Od3rY0bA}

\medskip
Animated GIF
\end{column}

\end{columns}

\bigskip

\begin{itemize}

\item A PDF cannot play a video or a GIF, so the slide shows a still and \myhref{https://webpressive.github.io}{WebPressive} plays the media over it

\item \lstinline|\animstill{height}{poster}{media}{source}| links the still to the media through \lstinline|wpmedia:| and attaches the media file to the PDF

\item The orange badge opens the source video; other PDF viewers show the still

\item Clip: CSU Systems Engineering Friday Talk (Nov.~3, 2023) \urlyoutube{N17Od3rY0bA}

\end{itemize}
\annotation{
This slide demonstrates the animstill command for animated media in WebPressive. The same clip is shown twice: on the left as an MP4 video and on the right as an animated GIF. The PDF carries a still of each, and WebPressive plays the media file over the still at its exact position. The animstill command adds a wpmedia link to the still, attaches the media file to the PDF so the PDF alone carries it, and draws an orange play badge on the still's bottom-left corner that links to the source video.
}
\end{myslidefragile}
